\chapter*{Prefacio}
\addcontentsline{toc}{chapter}{Prefacio}

Este libro nace de una necesidad pedagógica concreta: explicar la poda alpha--beta con el nivel de detalle que normalmente se reserva a los algoritmos clásicos en textos como \emph{Introduction to Algorithms}. La poda alpha--beta suele presentarse como una optimización intuitiva de Minimax, pero esa presentación deja abiertas preguntas esenciales: por qué la poda no cambia la respuesta, qué significan exactamente las cotas, de dónde salen las desigualdades, qué hipótesis son necesarias y cómo se ve el algoritmo cuando se aplica a juegos reales.

La exposición no presupone conocimientos previos. Cada símbolo se introduce antes de usarse. Cada definición tiene una lectura informal. Cada prueba se separa en caso base, hipótesis inductiva y paso inductivo cuando corresponde. Cada ejemplo de juego se desarrolla a través de estados, sucesores, utilidad, retropropagación y poda.

\begin{summarybox}
El cuerpo principal evita el código. El objetivo de esta versión es construir primero la comprensión matemática y visual: árboles, valores, cotas, invariantes y demostraciones. El código puede añadirse después en un apéndice independiente, pero no forma parte de esta etapa.
\end{summarybox}

El libro se desarrollará por entregas. Esta primera entrega establece el marco formal y comienza los ejemplos gráficos con tres en raya, damas y ajedrez. Las entregas posteriores ampliarán los árboles, las pruebas del paper de Knuth y Moore, el análisis probabilístico, los casos extremos y las optimizaciones modernas.
